\documentclass[chinese]{tutordesk}

\studio[NE]{Northbridge English Studio}
\student{Mina Park}
\lesson{7}
\lessondate{26 Oct 2026}
\doctype{Handout}
\exam{IELTS Academic \textperiodcentered{} Writing Task 2}
\title{Writing a PEEL Body Paragraph}
\subtitle{How to turn a point into a paragraph an examiner can follow}
\level{Band 6.0 to 7.0}
\timing{60 min}
\focus{Task Response, Coherence and Cohesion}

\begin{document}
\maketitle

\begin{objectives}
  \item write a first sentence that makes one clear, arguable claim;
  \item develop that claim with reasoning and a specific example;
  \item finish with a link that answers the question, not a summary that repeats it.
\end{objectives}

\section{The four moves}

\noindent
\begin{tabularx}{\linewidth}{@{}>{\sffamily\bfseries}p{7mm}>{\raggedright}p{36mm}>{\raggedright\arraybackslash}X@{}}
  \headlabel{Tag} & \headlabel{Move} & \headlabel{What it does} \\
  \midrule[0.6pt]
  P  & Point   & States one claim that directly answers the question. \\
  Ex & Explain & Shows \emph{why} or \emph{how} the claim is true: the reasoning. \\
  Eg & Example & Makes the reasoning concrete: a case, a scenario, a consequence. \\
  L  & Link    & Shows how the paragraph answers the question as it is worded. \\
  \bottomrule[0.4pt]
\end{tabularx}

\begin{keypoint}[Explain before you exemplify]
Many Band 6 paragraphs jump from the point straight to an example, so the reader has
to guess why the example matters. The Explain sentence is where your reasoning
becomes visible, and it is usually the move that lifts a paragraph into the Band 7 range.
\end{keypoint}

\section{One question, two paragraphs}

\noindent\emph{In many cities, public libraries are being replaced by online information
services. Do the advantages of this development outweigh the disadvantages?}

\begin{compare}
\begin{example}[Effective \textperiodcentered{} Band 7 range]
\seg{P}The strongest argument against closing libraries is that they offer free,
supervised space, which no website can provide.
\seg{Ex}For students in crowded homes, a quiet, heated room with reliable internet is
often the difference between studying and not studying at all.
\seg{Eg}A teenager revising for exams in a flat shared with four relatives, for instance,
may find the local library is the only place where sustained concentration is possible.
\seg{L}The loss of physical libraries would therefore hit the least advantaged hardest,
and this outweighs the convenience of online services.
\end{example}
\begin{nonexample}[Less effective \textperiodcentered{} Band 5 to 6]
Libraries are very important for people. Many people use libraries every day.
For example, my friend goes to the library at the weekend. Also, libraries have many
books and computers and they are free. In conclusion, libraries should not be replaced
by the internet.
\end{nonexample}
\end{compare}

\subsection{Why the second version scores lower}
\begin{itemize}
  \item The first sentence is a \gloss{truism}{不言自明的道理}: nobody would disagree, so there is nothing to argue.
  \item There is no Explain move. The example about a friend is personal and proves nothing about cities.
  \item “In conclusion” closes the essay, not a paragraph, and the last sentence never mentions advantages or disadvantages.
\end{itemize}

\begin{teachernote}
Ask them to label the weak paragraph with P, Ex, Eg and L before they look at the model.
Most students discover that two of the four moves are missing. Elicit the fix orally
before any writing.
\end{teachernote}

\section{Chunks for the Explain and Link moves}
\begin{vocab}
  \chunk{this is largely because}{claim, then \emph{this is largely because} + clause}%
    {Rents have risen sharply. This is largely because new housing has not kept pace with demand.}
  \chunk{which means that}{adds a consequence to the previous clause}%
    {Online courses can be paused, which means that working adults can study at their own pace.}
  \chunk{a case in point is}{introduces one specific example}%
    {A case in point is the bus lane scheme, which cut journey times within a year.}
  \chunk[/ˌdɪs.prəˈpɔː.ʃən.ət.li/]{hit someone disproportionately hard}{verb + adverb; who suffers most}%
    {Price rises hit low-income families disproportionately hard.}
  \chunk{outweigh the benefits of}{subject + \emph{outweigh} + noun phrase; links back to the question}%
    {These costs clearly outweigh the benefits of a shorter school week.}
\end{vocab}

\section{Self-test}

\begin{question}
Which sentence works best as the Point for this question:
\emph{Should governments tax sugary drinks?}
\begin{choices}
  \choice Sugary drinks are popular all over the world.
  \choice* A tax on sugary drinks would cut consumption in the group most at risk: children.
  \choice Many people have different opinions about taxes.
  \choice Sugar is bad for your health and your teeth.
\end{choices}
\why{B is the only option that makes a claim someone could disagree with. A and C are truisms; D does not mention tax.}
\end{question}

\begin{question}
Write an Explain sentence for this point: \emph{Remote work benefits employers as well as employees.}
\answerlines{2}
\begin{solution}
Model: \emph{This is largely because firms no longer need to rent large offices, which means
that money saved on property can be spent on staff and training.}
\end{solution}
\end{question}

\begin{selfcheck}[Check your own paragraph]
  \item My first sentence makes one claim that answers the question.
  \item My Explain sentence says why or how, rather than adding another fact.
  \item My example is specific: a person, a place, a scenario or a consequence.
  \item My last sentence uses the key words of the question.
\end{selfcheck}

\printanswers
\end{document}
